\section{Finite Element Discretization}

In this section, we present a finite element approximation of the optimality system derived in the 
previous sections.

\subsection{Spatial Discretization}

Let \(\mathcal T_h\) be a conforming triangulation of the domain \(D\) with mesh size \(h\).

We introduce the finite element space

\[
V_h
=
\left\{
v_h\in C^0(\overline D)
;:;
v_h|_K \in \mathbb P_1(K),
\quad
\forall K\in\mathcal T_h,
\quad
v_h=0 \text{ on } \partial D
\right\}.
\]

Let

\[
{\varphi_i}_{i=1}^{N}
\]

be a basis of \(V_h\).

The discrete state, adjoint, and control variables are approximated by

\[
y_h(x,t)
=
\sum_{i=1}^{N}
Y_i(t)\varphi_i(x),
\]

\[
p_h(x,t)
=
\sum_{i=1}^{N}
P_i(t)\varphi_i(x),
\]

and

\[
u_h(x,t)
=
\sum_{i=1}^{N}
U_i(t)\varphi_i(x).
\]

\subsection{Discrete Variational Formulation}

The finite element approximation of the state equation reads:

Find \(y_h(t)\in V_h\) such that

\begin{equation}
\left(
\frac{\partial y_h}{\partial t},
v_h
\right)
+
a(y_h,v_h)
=
(u_h,v_h),
\qquad
\forall v_h\in V_h.
\label{eq:state_discrete}
\end{equation}

Similarly, the adjoint problem becomes:

Find \(p_h(t)\in V_h\) such that

\begin{equation}
\left(
-\frac{\partial p_h}{\partial t},
v_h
\right)
+
a^\ast(p_h,v_h)
=
(y_h-y_d,v_h),
\qquad
\forall v_h\in V_h.
\label{eq:adjoint_discrete}
\end{equation}

\subsection{Finite Element Matrices}

Choosing the basis functions as test functions leads to the matrix formulation.

The mass matrix is defined by

\begin{equation}
M_{ij}
=
\int_D
\varphi_i \varphi_j
\,dx.
\label{eq:mass_matrix}
\end{equation}

The diffusion matrix is

\begin{equation}
K_{ij}
=
\kappa
\int_D
\nabla\varphi_i\cdot\nabla\varphi_j
\,dx.
\label{eq:stiffness_matrix}
\end{equation}

The advection matrix is

\begin{equation}
C_{ij}
=
\int_D
(\mathbf b\cdot\nabla\varphi_j)
\varphi_i
\,dx.
\label{eq:advection_matrix}
\end{equation}

The semi-discrete state system can therefore be written as

\begin{equation}
M\frac{dY}{dt}
+
(K+C)Y
=
MU.
\label{eq:semi_discrete_state}
\end{equation}

Similarly, the semi-discrete adjoint system is

\begin{equation}
-M\frac{dP}{dt}
+
(K+C^\ast)P
=
M(Y-Y_d).
\label{eq:semi_discrete_adjoint}
\end{equation}

\subsection{Time Discretization}

Let

\[
0=t_0<t_1<\cdots<t_N=T
\]

be a uniform partition of the time interval with time step

\[
\Delta t
=
\frac{T}{N}.
\]

Using the backward Euler scheme, the state equation becomes

\begin{equation}
\left(
\frac{M}{\Delta t}
+
K+C
\right)
Y^{n+1}
=
\frac{M}{\Delta t}Y^n
+
MU^{n+1}.
\label{eq:euler_state}
\end{equation}

The adjoint equation is solved backward in time:

\begin{equation}
\left(
\frac{M}{\Delta t}
+
K+C^\ast
\right)
P^{n}
=
\frac{M}{\Delta t}P^{n+1}
+
M(Y^n-Y_d^n).
\label{eq:euler_adjoint}
\end{equation}

\subsection{Discrete Optimality Condition}

The discrete counterpart of the optimality condition is

\begin{equation}
U^n
=
\frac{1}{\alpha}
P^n.
\label{eq:discrete_control}
\end{equation}

Consequently, the optimality system can be solved iteratively through the following procedure:

\begin{enumerate}
\item Initialize the control \(U^0\).
\item Solve the state equation forward in time.
\item Solve the adjoint equation backward in time.
\item Compute the gradient
\[
G^n
=
\alpha U^n-P^n.
\]
\item Update the control using a gradient-based algorithm.
\item Repeat until convergence.
\end{enumerate}

\subsection{Implementation in FEniCSx}

The finite element formulation presented above is well suited for implementation in FEniCSx.

The numerical implementation involves:

\begin{itemize}
\item mesh generation;
\item construction of the finite element space \(V_h\);
\item assembly of the mass, diffusion, and advection matrices;
\item solution of the forward state equation;
\item solution of the backward adjoint equation;
\item gradient computation and control update.
\end{itemize}

The resulting framework enables the numerical approximation of optimal control problems governed 
by stochastic advection--diffusion equations and serves as the basis for the numerical experiments 
presented in the next section.
